\section{紧空间到局部紧空间的嵌入}

\begin{theorem}{单点紧化}{one-point-conpactification}
	设$X$是局部紧空间.
	\begin{enumerate}
		\item 存在紧空间$X^{\prime}$,$\omega\in X^{\prime}$和$f\in\Isom_{\mathsf{Top}}\left(X,X^{\prime}\setminus\left\{\omega\right\}\right)$.
		\item 若存在紧空间$Y^{\prime}$,$\omega^{\prime}\in Y^{\prime}$和$f_{1}\in\Isom_{\mathsf{Top}}\left(X,Y^{\prime}\setminus\left\{\omega^{\prime}\right\}\right)$,则存在唯一的$g\in\Hom_{\mathsf{Top}}\left(X^{\prime},Y^{\prime}\right)$使$f_{1}=g\circ f$.
	\end{enumerate}
\end{theorem}

\begin{definition}{无穷原点,单点紧化}{}
	沿用定理(\ref{thm:one-point-conpactification})的设定.我们称$\omega$是$X^{\prime}$的无穷远点,$X^{\prime}$是$X$的单点紧化.
\end{definition}

\begin{remark}{}{}
	如果$X$是紧空间,则$\omega$是$X^{\prime}$的孤立点,并且$X^{\prime}=X\sqcup\left\{\omega\right\}$.
\end{remark}

\begin{example}{平面的单点紧化与立体投影}{}
	$\R^{2}$的单点紧化同胚于$\R^{3}$中的单位球面$\left\{\left(x_{1},x_{2},x_{3}\right)\in\R^{3}\middle|x_{1}^{2}+x_{2}^{2}+x_{3}^{2}=1\right\}$.具体的同胚定义如下:
	\begin{itemize}
		\item 把无穷远点$\omega$映射到$\left(0,0,1\right)$.
		\item 把$\R^{2}$的任意一点$\left(x_{1},x_{2}\right)$映射到链接$\left(0,0,1\right)$和$\left(x_{1},x_{2},0\right)$的直线与该球面的唯一交点.
	\end{itemize}
\end{example}










